% !TeX root = ../../Book2.tex
\chapter{一阶偏微分方程与特征线法}
\label{b2:ch:07}
\BookChapterNumber{b2:ch:07}
\begin{chapterabstract}
输运方程沿运动轨道记录函数值的变化。常速度给出平移，变速度则由常微分方程的流确定运动。本章使用\chapref{b2:ch:ode}已经证明的存在唯一性与一阶光滑依赖，建立含源项与零阶项的表示公式。支集传播描述信息从哪里到达，Jacobi 行列式描述体积怎样变化，两者共同给出密度的质量输运。选读部分讨论拟线性方程和 Hamilton--Jacobi 方程的局部光滑解，并以 Burgers 方程说明特征投影失去可逆性时会发生什么。
\end{chapterabstract}

\begin{learninggoals}
\begin{itemize}
\item\label{b2:item:07-goal-lines}能从一阶微分算子找出运动方向，沿特征线求解含源项与吸收项的初值问题。
\item\label{b2:item:07-goal-flow}能为特征轨道核对局部流的适用条件，利用唯一性与可微依赖构造非特征坐标，并区分局部回溯和全时间回溯。
\item\label{b2:item:07-goal-density}能用回溯轨道判断支集传播，并用 Jacobi 公式区分输运标量与守恒密度。
\item\label{b2:item:07-goal-nonlinear}选读目标：在\secref{b2:sec:0704}的局部可逆条件下构造拟线性解，辨认梯度爆破与振幅爆破的区别，理解弱守恒关系的有限作用；只学习经典线性方程时可暂缓这一目标。
\end{itemize}
\end{learninggoals}

一团染料被介质携带时，它的边界随粒子轨道移动。若介质同时膨胀，染料浓度还会降低；仅把初始图形平移或变形，未必保存总质量。要由瞬时速度计算以后每一点的数值，需要知道轨道怎样返回初值面、相邻轨道是否可区分，以及体积怎样变化。特征线法逐个解决这些问题，并把一阶偏微分方程转为沿轨道的常微分方程。

\ProseParagraphBreak

本章使用\chapref{b2:ch:01}的链式法则、逆函数定理、换元及积分求导规则，以及\chapref{b2:ch:ode}已写部分的局部存在唯一性、线性增长延拓和一阶光滑依赖。\secref{b2:sec:0701}从便于计算的直线模型讲起；\secref{b2:sec:0702}沿特征轨道建立表示公式，并用非特征数据曲线构造局部坐标；\secref{b2:sec:0703}以特征流求解时变输运方程，并研究质量守恒。只学习经典线性方程的读者可以略过\secref{b2:sec:0704}，后续普通章节不需要该节的非线性结论。

\ProseParagraphBreak

Walter A. Strauss 的《Partial Differential Equations: An Introduction》\cite[Section 1.2]{Strauss2008Partial}由直线上的求解过渡到特征曲线。局部流的存在唯一性与一阶光滑依赖见\chapref{b2:ch:ode}；本章进一步计算沿流的函数值及体积变化。
% 实读：Strauss 第二版 Wiley 官方样章印刷6--9页；常微分方程来源说明随已有证明移至第2章。
% 第7--8章来源实读范围与改题出处见参考文献原文件/来源记录.md。

% !TeX root = ../../Book2.tex
\section{一阶线性方程}
\label{b2:sec:0701}

给定某一时刻的浓度分布，若介质只作平移而没有扩散，后来的图形应当保持原来的形状。我们先用方程描述这种平移，再沿直线求出解，由此介绍特征线法。本章把空间变量写在前面，时间变量写在后面，即 \(u=u(x,t)\)。向量场表示速度，函数值既可能表示随粒子携带的量，也可能表示单位体积中的密度；这两种含义将在\secref{b2:sec:0703}区分。

\subsection{输运方程}
\label{b2:subsec:070101}

设 \(d\geq1\) 为整数，\(b\in\mathbb R^d\) 为固定速度。经过初始位置 \(a\) 的粒子在时刻 \(t\) 位于 \(a+tb\)。对 \(u\in C^1(\mathbb R^d\times I;\mathbb C)\)，其中 \(I\) 是时间开区间，若沿这一运动携带的数值不变，则链式法则给出
\[
 \frac{\mathrm d}{\mathrm dt}u(a+tb,t)
 =u_t(a+tb,t)+b\cdot\nabla_xu(a+tb,t)=0.
\]
这产生\term[shuyun-fangcheng]{输运方程}[transport equation]
\begin{equation}\label{b2:eq:07-constant-transport}
 u_t+b\cdot\nabla_xu=0.
\end{equation}
其中 \(u_t\) 衡量固定观察点的变化，\(b\cdot\nabla_xu\) 衡量移动造成的变化。二者可以分别非零而恰好相消。例如一维函数 \(u(x,t)=\sin(x-bt)\) 的每个波峰向右或向左平移，方向取决于 \(b\) 的符号。

\ProseParagraphBreak

这里讨论的经典解属于 \(C^1\)，因此上述链式法则有逐点意义。对于不光滑的解，本卷\chapref{b2:ch:05}的分布语言把方程改写为测试函数恒等式。

\subsection{常系数情形}
\label{b2:subsec:070102}

\begin{proposition}\label{b2:prop:07-translation-solution}
设 \(d\geq1\) 为整数，\(b\in\mathbb R^d\) 为常向量，\(I\) 为包含 \(0\) 且有非空内部的时间区间，\(u\in C^1(\mathbb R^d\times I;\mathbb C)\)；若 \(I\) 含端点，\(C^1\) 按能局部延伸到邻域理解。函数 \(u\) 满足 \(u_t+b\cdot\nabla u=0\)，当且仅当存在 \(g\in C^1(\mathbb R^d;\mathbb C)\)，使
\[
 u(x,t)=g(x-tb).
\]
其中 \(g(x)=u(x,0)\)，因而初值唯一确定解。
\end{proposition}
\begin{proof}
固定 \((x,t)\)，取 \(a=x-tb\)。函数 \(s\mapsto u(a+sb,s)\) 的导数为零，从 \(0\) 积分到 \(t\) 得 \(u(x,t)=u(a,0)\)。反过来，对所给公式求导，有 \(u_t=-b\cdot\nabla g(x-tb)\)、\(\nabla_xu=\nabla g(x-tb)\)，故方程成立。时间为负时，积分的方向反转，论证不变。
\end{proof}

平移不会提高正则性：初值在哪里存在折角，后来就在相应平移点保留折角。对于任意 \(1\leq p\leq\infty\)，若初值还属于 \(L^p\)，平移不变性给出 \(\|u(\cdot,t)\|_p=\|g\|_p\)。这一恒等式没有衰减因子，与以后热方程的平滑效应不同。

\ProseParagraphBreak

不把某个变量预先叫作时间时，同样的方法适用于 \(\alpha u_x+\beta u_y=0\)，其中 \((\alpha,\beta)\in\mathbb R^2\setminus\{(0,0)\}\)、\(u\in C^1(\mathbb R^2;\mathbb C)\)。取
\[
 r=\alpha x+\beta y,\quad s=\beta x-\alpha y,
\]
则变换的行列式为 \(-\alpha^2-\beta^2\ne0\)，算子化为 \((\alpha^2+\beta^2)\partial_r\)。所以在整个平面上，解恰为 \(F(\beta x-\alpha y)\)，其中 \(F\in C^1(\mathbb R;\mathbb C)\)。若区域与一条直线的交集不连通，各个连通片上的常数未必由同一个全局函数 \(F\) 描述；局部坐标公式不能自动升级为任意区域上的全局表示。

\subsection{非齐次方程}
\label{b2:subsec:070103}

源项使粒子所携带的值随时间累积，实数零阶系数造成吸收或增长。为与 Fourier 分析中的复值函数兼容，下面允许常数 \(\lambda\in\mathbb C\)，并明确让初值、源项和解都取复值。因子 \(e^{-\lambda t}\) 的模为 \(e^{-(\operatorname{Re}\lambda)t}\)，虚部还会改变相位。

\begin{theorem}\label{b2:thm:07-constant-source}
设 \(d\geq1\) 为整数，\(T>0\)，\(b\in\mathbb R^d\)、\(\lambda\in\mathbb C\) 为常数，\(g\in C^1(\mathbb R^d;\mathbb C)\)、\(f\in C^1(\mathbb R^d\times[0,T];\mathbb C)\)。问题
\[
 u_t+b\cdot\nabla u+\lambda u=f,\quad u(x,0)=g(x)
\]
在 \(\mathbb R^d\times[0,T]\) 上有唯一的复值 \(C^1\) 解，且
\begin{equation}\label{b2:eq:07-constant-source}
 u(x,t)=e^{-\lambda t}\cdot g(x-tb)
 +\int_0^t e^{-\lambda(t-s)}\cdot f\bigl(x-(t-s)b,s\bigr)\,\mathrm ds.
\end{equation}
闭时间区间上的光滑性均理解为能局部延伸到邻域，端点导数也可取相应单侧值。
\end{theorem}
\begin{proof}
沿 \(x=a+tb\) 令 \(v(t)=u(a+tb,t)\)，方程成为
\(v'+\lambda v=f(a+tb,t)\)。乘以 \(e^{\lambda t}\) 并从 \(0\) 积分，得到
\[
 v(t)=e^{-\lambda t}\cdot g(a)
 +\int_0^t e^{-\lambda(t-s)}\cdot f(a+sb,s)\,\mathrm ds.
\]
令 \(a=x-tb\) 即为表示公式，故任何经典解只能取此值。公式右侧在每个紧时空集合上的被积函数及一阶导数都有共同界；\cref{b2:thm:01-parameter-differentiation}允许对空间参数在积分内求导，\cref{b2:prop:01-leibniz-integral}还处理时间上限的变化，因而右侧属于 \(C^1\)。把它沿每条直线拉回，所得函数满足刚才的常微分方程；链式法则再给出原方程。取 \(t=0\) 得到初值。
\end{proof}

例如在一维取 \(b=1\)、\(\lambda=0\)、\(g=0\)、\(f(x,t)=x\)，有
\[
 u(x,t)=\int_0^t(x-t+s)\,\mathrm ds=tx-\frac{t^2}{2}.
\]
这里源项须在粒子过去所在的位置取值；若误把积分写成 \(\int_0^tf(x,s)\,\mathrm ds\)，会得到不满足方程的 \(tx\)。

复值范围也有直接的例子：取 \(\lambda=\mathrm i\)、\(g\equiv1\)、\(f\equiv0\)，就得到 \(u(x,t)=e^{-\mathrm it}\)。其模保持为 \(1\)，函数值却一般不是实数。若系数 \(\lambda\)、初值与源项都为实值，表示公式自然给出实值解。

\subsection{初值问题}
\label{b2:subsec:070104}

初值面 \(t=0\) 被每条轨道恰好穿过一次，因此可以给每条轨道一个独立数值。若改在与轨道相切的曲线上给值，就不能仅凭一般的特征坐标构造保证存在性与唯一性。以 \(u_x+u_y=0\) 为例，沿直线 \(y=x\) 解必须恒定；指定 \(u(s,s)=s\) 没有解，而指定 \(u(s,s)=0\) 则允许所有满足 \(F(0)=0\) 的解 \(u(x,y)=F(x-y)\)。方程与数据之间的几何位置是初值问题的一部分。

\ProseParagraphBreak

常系数公式还直接给出数据稳定性。固定整数 \(d\geq1\)、\(T>0\)、\(b\in\mathbb R^d\) 与 \(\lambda\in\mathbb C\)。若 \(u,\widetilde u\in C^1(\mathbb R^d\times[0,T];\mathbb C)\) 分别满足同一算子的方程 \(u_t+b\cdot\nabla u+\lambda u=f\) 与 \(\widetilde u_t+b\cdot\nabla\widetilde u+\lambda\widetilde u=\widetilde f\)，初值分别为 \(g,\widetilde g\)，则对 \(0\leq t\leq T\)，只要右侧范数有限，就有
\[
 \|u(\cdot,t)-\widetilde u(\cdot,t)\|_\infty
 \leq e^{|\lambda|t}\left(\|g-\widetilde g\|_\infty
 +\int_0^t\|f(\cdot,s)-\widetilde f(\cdot,s)\|_\infty\,\mathrm ds\right).
\]
证明只需在\eqref{b2:eq:07-constant-source}中相减并取绝对值。下一节把直线改成由速度场决定的曲线；为使这种替换合法，先要证明曲线存在、彼此不发生碰撞，并且随初始点可微地变化。

% !TeX root = ../../Book2.tex
\section{特征曲线}
\label{b2:sec:0702}

变系数速度场已经给出粒子在每一点的瞬时运动方向，却没有直接给出整条轨道。先解常微分方程，再用轨道作坐标，是本节的构造顺序。以下 \(U\subseteq\mathbb R^d\) 为开集，\(b\in C^1(U;\mathbb R^d)\)；本节先处理不显含时间的速度场。

\subsection{特征常微分方程}
\label{b2:subsec:070201}

\begin{definition}\label{b2:def:07-characteristic}
设 \(d\) 为正整数，\(U\subset\mathbb R^d\) 为开集，\(b:U\to\mathbb R^d\) 连续，\(I\subset\mathbb R\) 为区间。若 \(X:I\rightarrow U\) 属于 \(C^1\)，并满足
\[
 X'(t)=b\bigl(X(t)\bigr),
\]
则称它为速度场 \(b\) 的\term[tezheng-quxian]{特征曲线}[characteristic curve]，也称特征线。其常微分方程称为\term[tezheng-changweifen-fangcheng]{特征常微分方程}[characteristic ordinary differential equation]。
\end{definition}

沿这样的曲线，复合函数的导数只涉及 \(u_t+b\cdot\nabla u\)。对本节采用的 \(C^1\) 速度场，\cref{b2:thm:07-local-ode}保证经过每个初始点的轨道局部存在且唯一。将其最大解记为 \(\Phi_t(a)\)；流的复合、逆关系和对初始点的 \(C^1\) 依赖已经在\secref{b2:sec:ode-smooth-flow}证明。下面利用这些结论，把沿轨道求出的函数还原为偏微分方程的经典解。

\subsection{沿特征线的表示公式}
\label{b2:subsec:070203}

对自治输运问题 \(u_t+b(x)\cdot\nabla u=0\)，沿轨道有
\[
 \frac{\mathrm d}{\mathrm dt}u\bigl(\Phi_t(a),t\bigr)=0.
\]
因此只要轨道可从时刻 \(t\) 返回 \(0\)，解就由
\begin{equation}\label{b2:eq:07-autonomous-representation}
 u(x,t)=g\bigl(\Phi_{-t}(x)\bigr)
\end{equation}
给出。这里的“返回”既包括时间存在范围，也包括轨道始终位于系数定义域；若在有界区域中轨道先碰到边界，单靠内部初值不足以继续决定数值。

\ProseParagraphBreak

例如平面速度 \(b(x,y)=(-y,x)\) 给出
\[
 \Phi_t(a_1,a_2)=(a_1\cos t-a_2\sin t,\ a_1\sin t+a_2\cos t).
\]
解是初始图形的旋转，\(u(x,y,t)=g(x\cos t+y\sin t,-x\sin t+y\cos t)\)。另一例是 \(b(x)=x\)，其流为 \(\Phi_t(a)=e^ta\)，故 \(u(x,t)=g(e^{-t}x)\)。后者虽然函数值不变，体积却按 \(e^{dt}\) 扩张；把 \(u\) 的积分解释成守恒质量就会出错。

\subsection{全局流与非特征初值问题}
\label{b2:subsec:070204}

\begin{proposition}\label{b2:prop:07-global-flow}
设 \(d\) 为正整数，\(b\in C^1(\mathbb R^d;\mathbb R^d)\)，且存在 \(C\geq0\) 使对所有 \(x\in\mathbb R^d\) 都有
\( |b(x)|\leq C(1+|x|)\)，记 \(\Phi_t\) 为 \(b\) 的自治流。对 \(g\in C^1(\mathbb R^d;\mathbb C)\)，\(u(x,t)=g(\Phi_{-t}(x))\) 给出 \(u_t+b(x)\cdot\nabla u=0\)、\(u(x,0)=g(x)\) 在 \(\mathbb R^d\times\mathbb R\) 上的唯一复值 \(C^1\) 解。
\end{proposition}
\begin{proof}
由\cref{b2:prop:ode-linear-growth}，所有轨道在整个实轴上存在；\cref{b2:thm:07-flow-derivative}给出全局流的 \(C^1\) 性。因此表示公式给出 \(C^1\) 函数。沿特征线该函数等于 \(g(a)\)，故满足方程。任意另一经典解沿同一轨道也恒等于 \(g(a)\)，因此唯一。
\end{proof}

流也给出一般数据曲线的局部求解。把平面上的向量场写成 \(\beta=(\beta_1,\beta_2)\)，给定参数曲线 \(\gamma(s)\) 后，重要的是它是否横穿轨道，而不是曲线能否写成水平直线。

\begin{proposition}[非特征数据曲线]\label{b2:prop:07-noncharacteristic-data}
设 \(U\subset\mathbb R^2\) 为开集，\(\beta\in C^1(U;\mathbb R^2)\)、\(c,f\in C^1(U;\mathbb R)\)，\(I\subset\mathbb R\) 为开区间，\(\gamma\in C^1(I;U)\)、\(g\in C^1(I;\mathbb R)\)，\(s_0\in I\)，且在 \(s=s_0\) 有
\[
 \det\bigl(\gamma'(s_0),\beta(\gamma(s_0))\bigr)\ne0.
\]
则在 \(\gamma(s_0)\) 附近，问题 \(\beta\cdot\nabla u+cu=f\)、\(u(\gamma(s))=g(s)\) 有唯一的局部实值 \(C^1\) 解；这里的曲线数据只对 \(s_0\) 附近的参数规定。局部唯一性指任意两个满足同一曲线数据的局部解，在适当缩小的共同邻域内相等。
\end{proposition}
\begin{proof}
记 \(\Phi_t\) 为 \(\beta\) 的局部流，定义 \(\Psi(s,t)=\Phi_t(\gamma(s))\)。流对时间和初始点为 \(C^1\)，故 \(\Psi\) 也为 \(C^1\)。在 \((s_0,0)\)，它的两列导数正是 \(\gamma'(s_0)\)、\(\beta(\gamma(s_0))\)，故\cref{b2:thm:01-inverse}给出局部 \(C^1\) 逆坐标。缩小到参数矩形 \(J\times(-\tau,\tau)\)，其中 \(J\subset I\) 是含 \(s_0\) 的开区间、\(\tau>0\)，令 \(W=\Psi(J\times(-\tau,\tau))\subset U\)，则 \((s,t)=\Psi^{-1}(x)\) 在 \(W\) 上有定义。沿每个固定 \(s\) 解标量方程
\[
 \partial_tV(s,t)+c(\Psi(s,t))V(s,t)=f(\Psi(s,t)),
 \quad V(s,0)=g(s).
\]
具体地，记
\[
 C(s,t)=\int_0^t c(\Psi(s,r))\,\mathrm dr.
\]
积分因子公式给出唯一的
\[
 V(s,t)=e^{-C(s,t)}\left(g(s)+
       \int_0^t e^{C(s,r)}f(\Psi(s,r))\,\mathrm dr\right).
\]
对参数矩形中的每个紧子矩形，被积函数及一阶导数有共同界，\secref{b2:sec:0103}的积分求导规则于是给出 \(V\in C^1\)。令 \(u=V\circ\Psi^{-1}\)。由于 \(\partial_t\Psi=\beta(\Psi)\)，链式法则给出方程；\(t=0\) 给出曲线数据。任何另一经典解沿同一流也满足这个标量方程，所以在缩小后的坐标邻域内相同。
\end{proof}

设 \(U\subset\mathbb R^2\) 开，\(\beta\in C^1(U;\mathbb R^2)\)，\(\gamma:I\to U\) 为 \(C^1\) 曲线，其中 \(I\) 为开区间。若对每个 \(s\in I\) 都有 \(\det(\gamma'(s),\beta(\gamma(s)))\ne0\)，则称它为\term[fei-tezheng-quxian]{非特征曲线}[noncharacteristic curve]。例如对 \(u_x+u_y=f\) 和 \(\gamma(s)=(s,\kappa s)\)，其中 \(\kappa\in\mathbb R\)，行列式为 \(1-\kappa\)，因此 \(\kappa\ne1\) 时局部数据可独立指定；\(\kappa=1\) 时曲线沿轨道，数据必须满足 \(g'(s)=f(s,s)\)。上述非退化条件使曲线切向和流方向共同组成局部坐标。

\ProseParagraphBreak

\Cref{b2:fig:07-noncharacteristic-data}取同一个常向量场 \(\beta=(1,0)\)，其流为 \(\Phi_t(a_1,a_2)=(a_1+t,a_2)\)。左图取 \(\gamma(s)=(0,s)\)，则 \(\Psi(s,t)=(t,s)\)：数据曲线横穿每条附近的水平轨道，两个参数确实给出二维坐标。右图取 \(\gamma(s)=(s,0)\)，则 \(\Psi(s,t)=(s+t,0)\)：初值参数和流时间都只改变同一条轨道上的位置，不能用它们覆盖二维邻域。

% 本书原创：同一尺度绘出 beta=(1,0) 的精确水平流和两条数据曲线。
\begin{figure}[htbp]
\centering
\begin{tikzpicture}[x=1cm,y=1cm,>=stealth,every node/.style={font=\small},line width=.65pt]
\begin{scope}
\draw[->,gray!60] (-2.2,0)--(2.35,0) node[right] {\(x\)};
\draw[->,gray!60] (0,-1.7)--(0,1.85) node[above] {\(y\)};
\foreach \eta in {-1,0,1}{
  \draw[->,AnalysisGray] (-2.05,\eta)--(2.02,\eta);
}
\draw[->,AnalysisGray,very thick] (.4,1)--(1.4,1) node[midway,above=4pt] {\(\beta=(1,0)\)};
\draw[->,AnalysisBlue,very thick] (0,-1.5)--(0,1.5);
\node[AnalysisBlue,left] at (-.08,1.35) {\(\gamma\)};
\fill[AnalysisBlue] (0,0) circle(1.4pt);
\node at (0,2.6) {横截：\(\gamma(s)=(0,s)\)};
\node at (0,-2.15) {\(\det(\gamma',\beta)=-1\ne0\)};
\end{scope}
\begin{scope}[xshift=6.4cm]
\draw[->,gray!60] (-2.2,0)--(2.35,0) node[right] {\(x\)};
\draw[->,gray!60] (0,-1.7)--(0,1.85) node[above] {\(y\)};
\foreach \eta in {-1,0,1}{
  \draw[->,AnalysisGray] (-2.05,\eta)--(2.02,\eta);
}
\draw[->,AnalysisGray,very thick] (.4,1)--(1.4,1) node[midway,above=4pt] {\(\beta=(1,0)\)};
\draw[->,AnalysisBlue,very thick] (-1.5,0)--(1.5,0);
\node[AnalysisBlue,above] at (-.75,.08) {\(\gamma\)};
\fill[AnalysisBlue] (0,0) circle(1.4pt);
\node at (0,2.6) {相切：\(\gamma(s)=(s,0)\)};
\node at (0,-2.15) {\(\det(\gamma',\beta)=0\)};
\end{scope}
\end{tikzpicture}
\caption[数据曲线与特征轨道的横截和相切]{数据曲线与特征轨道的横截和相切。两图均取 \(\beta=(1,0)\)，灰线为水平轨道，粗灰箭头为长度为 \(1\) 的速度向量，蓝线为数据曲线；两图使用同一坐标尺度。左图的数据曲线与每条附近轨道唯一相交，\(\Psi(s,t)=(t,s)\) 可作局部坐标。右图的数据曲线完全落在一条轨道上，\(\Psi(s,t)=(s+t,0)\) 退化；它仅展示退化可能造成的失败，不代表任意行列式为零的情形都无解或不唯一。}
\label{b2:fig:07-noncharacteristic-data}
\end{figure}


在右图的具体模型中，取方程 \(u_x=0\) 和曲线数据 \(u(s,0)=s\)，沿曲线求导会同时得到 \(u_x(s,0)=1\) 与 \(u_x(s,0)=0\)，所以局部无解。若改取常数数据 \(g\equiv c_0\)，则 \(u(x,y)=c_0+q(y)\) 对任意满足 \(q(0)=0\) 的实值 \(C^1\) 函数 \(q\) 都给出延拓。这个模型说明退化时可能出现的失败方式，并不把一般的行列式为零判作必然无解或必然不唯一；前一命题的局部构造依靠的是逆函数定理的非退化假设。

% !TeX root = ../../Book2.tex
\section{变系数输运方程}
\label{b2:sec:0703}

本节把时间依赖、源项和体积变化同时纳入输运。这里 \(d\) 为正整数，\(T>0\)。为避免轨道在有限时间逃离空间，先假定 \(b\in C^1(\mathbb R^d\times[0,T];\mathbb R^d)\)，并有 \(|b(x,t)|\leq C(1+|x|)\)、\(C\geq0\)。局部版本只需把讨论限制在能够返回初始面的轨道邻域中。

\subsection{时空特征线}
\label{b2:subsec:070301}

把从时刻 \(s\) 的位置 \(a\) 出发的轨道记为 \(X(t;s,a)\)，即
\[
 \partial_tX(t;s,a)=b\bigl(X(t;s,a),t\bigr),\quad X(s;s,a)=a.
\]
扩充空间中的向量场 \(B(x,t)=(b(x,t),1)\) 是自治的。对 \(B\) 应用\cref{b2:thm:07-local-ode}与\cref{b2:thm:07-flow-derivative}，再以空间线性增长估计延拓，便得到整个 \([0,T]\) 上的轨道及其 \(C^1\) 依赖。特别地，时间坐标严格按单位速度增加，故轨道与每个等时面恰相交一次。
\symidx{\(X(t;s,a)\)：从时刻 \(s\)、位置 \(a\) 出发的轨道}

\ProseParagraphBreak

唯一性给出的复合关系现在是
\[
 X(t;r,X(r;s,a))=X(t;s,a),\quad
 X(s;t,X(t;s,a))=a.
\]
中间时刻不能省略：一般没有仅由 \(t-s\) 决定的单参数流。固定终点 \((x,t)\) 后，曲线 \(r\mapsto X(r;t,x)\) 描述该粒子的过去，初始位置为 \(X(0;t,x)\)。例如 \(b(x,t)=t x\) 时，
\[
 X(t;s,a)=\exp\left(\frac{t^2-s^2}{2}\right)a;
\]
从 \(s=0\) 走一个单位时间与从 \(s=1\) 走一个单位时间，伸缩比例不同。

\begin{proposition}[速度扰动的轨道估计]\label{b2:prop:07-speed-stability}
设 \(d\) 为正整数，\(T>0\)，\(L,\delta\geq0\)，\(b,\widetilde b\in C^1(\mathbb R^d\times[0,T];\mathbb R^d)\) 满足 \(|b(x,t)|+|\widetilde b(x,t)|\leq C(1+|x|)\)（某个 \(C\geq0\)），从 \((a,s)\)、\(a\in\mathbb R^d\) 出发的轨道分别记为 \(X,\widetilde X\)，\(s,t\in[0,T]\)，并且
\[
 \sup_{x,t}\|D_xb(x,t)\|\leq L,
 \quad\sup_{x,t}|b(x,t)-\widetilde b(x,t)|\leq\delta.
\]
则从同一时刻、同一点出发的两条轨道满足
\[
 |X(t;s,a)-\widetilde X(t;s,a)|
 \leq\delta|t-s|e^{L|t-s|}.
\]
\end{proposition}
\begin{proof}
将两条轨道方程相减，在速度差中加减 \(b(\widetilde X,r)\)。中值估计和速度差界给出，若 \(t\geq s\)，
\[
 |X(t)-\widetilde X(t)|
 \leq L\int_s^t|X(r)-\widetilde X(r)|\,\mathrm{d}r+\delta(t-s).
\]
对一个固定终点 \(t\)，把右侧最后一项在整个 \([s,t]\) 上放大为常数 \(\delta(t-s)\)，Gronwall 给出所述界。反向时间对反向速度应用相同论证，所以 \(t<s\) 也成立。
\end{proof}

若同一初值 \(g\in C^1(\mathbb R^d;\mathbb C)\) 有全空间 Lipschitz 常数 \(M\)，分别以这两个速度场求解 \(u_t+b\cdot\nabla u=0\) 与 \(\widetilde u_t+\widetilde b\cdot\nabla\widetilde u=0\)，便有 \(\|u(\cdot,t)-\widetilde u(\cdot,t)\|_\infty\leq M\delta t e^{Lt}\)、\(0\leq t\leq T\)。这里没有源项或零阶项；对终点固定的两条回溯轨道使用命题，再将它们的初始位置代入 \(g\)。例如常速度增加 \(\varepsilon\) 时，相同图形的位置误差为 \(\varepsilon t\)。若初始图形有跳跃，平移再小也可能在上确界范数中产生大小为一的误差；速度稳定性和选择何种数据范数必须同时考虑。

\subsection{源项}
\label{b2:subsec:070302}

\begin{theorem}\label{b2:thm:07-variable-transport}
设 \(d\) 为正整数，\(T>0\)，\(b\in C^1(\mathbb R^d\times[0,T];\mathbb R^d)\)，且 \(|b(x,t)|\leq C(1+|x|)\) 对某个 \(C\geq0\) 成立。记 \(X(t;s,a)\) 为 \(\partial_tX=b(X,t)\)、\(X(s;s,a)=a\) 的轨道。设 \(c,f\in C^1(\mathbb R^d\times[0,T];\mathbb C)\)、\(g\in C^1(\mathbb R^d;\mathbb C)\)。初值问题
\[
 u_t+b\cdot\nabla u+c\cdot u=f,\quad u(x,0)=g(x)
\]
存在唯一的经典解 \(u\in C^1(\mathbb R^d\times[0,T];\mathbb C)\)。记
\[
 E(t,r;x)=\exp\left(-\int_r^t c(X(q;t,x),q)\,\mathrm dq\right),
\]
则解为
\begin{equation}\label{b2:eq:07-variable-transport}
 u(x,t)=E(t,0;x)\cdot g(X(0;t,x))
 +\int_0^tE(t,r;x)\cdot f(X(r;t,x),r)\,\mathrm dr.
\end{equation}
\end{theorem}
\begin{proof}
固定 \(a\)，令 \(v(t)=u(X(t;0,a),t)\)，则链式法则把原方程变成
\[
 v'(t)+c(X(t;0,a),t)\cdot v(t)=f(X(t;0,a),t),\quad v(0)=g(a).
\]
用积分因子求解，再令 \(a=X(0;t,x)\)，即得\eqref{b2:eq:07-variable-transport}。因每条轨道的标量方程唯一，原方程也至多有一个经典解。

反过来，所给公式在紧时空集合上由 \(C^1\) 函数的复合及有限区间积分组成。轨道的连续性使相应轨道集合紧致，因此被积函数及一阶导数均有局部共同界，积分求导合法。沿 \(X(t;0,a)\) 拉回该公式时，复合关系把每条历史轨道化为 \(X(r;0,a)\)，于是恢复刚才的积分因子公式，得到方程和初值。
\end{proof}

若 \(c\) 为实值且 \(c\geq0\)，因子 \(E\) 介于 \(0\) 与 \(1\) 之间。对于非负实值初值和源项，解始终非负；对于有界复值数据，
\[
 \|u(\cdot,t)\|_\infty\leq\|g\|_\infty
 +\int_0^t\|f(\cdot,r)\|_\infty\,\mathrm dr.
\]
对一般复值 \(c\)，若存在 \(C_0\geq0\) 使 \(\operatorname{Re}c\geq-C_0\)，则 \(|E(t,r;x)|\leq e^{C_0(t-r)}\)，只需在上述上界右端补上 \(e^{C_0t}\)。这类估计还可用于两组数据之差，给出固定速度和零阶系数下的稳定性；两种上界均要求右侧的数据范数积分有限。源项按轨道积分，而零阶项以指数积累，二者对解的作用不同。

\subsection{支集传播}
\label{b2:subsec:070303}

在齐次情形 \(f=0\) 中，指数因子永不为零，因而
\begin{equation}\label{b2:eq:07-support-flow}
 \operatorname{supp}u(\cdot,t)
 =X(t;0,\operatorname{supp}g).
\end{equation}
这是闭支集的等式：流是全空间的同胚，它把非零点集合映为非零点集合，也把闭包映为闭包。

\ProseParagraphBreak

有源项时，各时刻产生的数值继续沿轨道传播，但不同来源之间可能抵消。由表示公式只能断言
\[
 \operatorname{supp}u(\cdot,t)\subseteq
 \overline{X(t;0,\operatorname{supp}g)
 \cup\bigcup_{0\leq r\leq t}X(t;r,\operatorname{supp}f(\cdot,r))}.
\]
为核对这一包含，取右侧闭集之外一点及一个与它不交的小邻域。每个邻近终点的回溯轨道都避开初始支集与全部源项支集，表示公式的每一项因此都为零。闭包不能随意去掉，因为不同时间的支集之并一般不闭。

\begin{proposition}\label{b2:prop:07-finite-speed}
设 \(d\) 为正整数，\(T>0\)，固定速度场 \(b\in C^1(\mathbb R^d\times[0,T];\mathbb R^d)\) 和零阶系数 \(c\in C^1(\mathbb R^d\times[0,T];\mathbb C)\)。设 \(M\geq0\) 且 \(|b(x,t)|\leq M\)，记 \(X(t;s,a)\) 为 \(\partial_tX=b(X,t)\)、\(X(s;s,a)=a\) 的轨道。则对 \(s,t\in[0,T]\)、\(a\in\mathbb R^d\)，有 \(|X(t;s,a)-a|\leq M|t-s|\)。

对 \(j=1,2\)，取 \(g_j\in C^1(\mathbb R^d;\mathbb C)\)、\(f_j\in C^1(\mathbb R^d\times[0,T];\mathbb C)\)，令 \(u_j\in C^1(\mathbb R^d\times[0,T];\mathbb C)\) 为同一组系数下 \((u_j)_t+b\cdot\nabla u_j+cu_j=f_j\)、\(u_j(x,0)=g_j(x)\) 的唯一经典解。对固定的 \((x,t)\in\mathbb R^d\times[0,T]\)，若 \(g_1=g_2\) 在 \(\overline B_{Mt}(x)\) 上成立，且 \(f_1=f_2\) 在回溯圆锥
\[
 \{\,(y,r)\mid0\leq r\leq t,\ |y-x|\leq M(t-r)\,\}
\]
上成立，则 \(u_1(x,t)=u_2(x,t)\)。因而在固定 \(b,c\) 后，该点的解值仅依赖于初值球内的 \(g\) 和回溯圆锥内的 \(f\)。
\end{proposition}
\begin{proof}
积分轨道方程并取绝对值得到速度界。对每个回溯时刻 \(r\)，该界给出 \(|X(r;t,x)-x|\leq M(t-r)\)，特别是初始位置属于 \(\overline B_{Mt}(x)\)。两解使用同一个流，也使用由同一 \(c\) 决定的积分因子 \(E\)。在\eqref{b2:eq:07-variable-transport}中相减，得到
\[
 \begin{aligned}
 u_1(x,t)-u_2(x,t)
 &=E(t,0;x)(g_1-g_2)(X(0;t,x))\\
 &\quad+\int_0^t E(t,r;x)(f_1-f_2)(X(r;t,x),r)\,\mathrm dr=0.
 \end{aligned}
\]
初值项和源项分别因球内和圆锥内的数据相同而为零。
\end{proof}

实际依赖集通常比这个圆锥小得多：标量输运只读取一条回溯曲线。圆锥的作用是给出不必显式求流的统一上界。线性增长的速度仍保证有限时间的轨道存在，但没有一个对所有起点通用的有限传播速度；例如 \(b(x)=x\) 在远离原点处可以运动得任意快。

\subsection{Jacobi 行列式与质量输运初步}
\label{b2:subsec:070304}

记
\[
 J(t;s,a)\coloneqq\det D_aX(t;s,a).
\]
初始时刻 \(J(s;s,a)=1\)。正、逆流的可微性保证导数矩阵一直可逆，所以行列式不能变号。下面的公式进一步说明其大小。
\symidx{\(J(t;s,a)\)：流对初始位置的 Jacobi 行列式}

\begin{proposition}[Liouville 公式]\label{b2:prop:07-jacobian}
设 \(d\) 为正整数，\(T>0\)，\(b\in C^1(\mathbb R^d\times[0,T];\mathbb R^d)\)，且 \(|b(x,t)|\leq C(1+|x|)\) 对某个 \(C\geq0\) 成立。记 \(X(t;s,a)\) 为 \(\partial_tX=b(X,t)\)、\(X(s;s,a)=a\) 的轨道。对 \(s,t\in[0,T]\)、\(a\in\mathbb R^d\)，令 \(J(t;s,a)=\det D_aX(t;s,a)\)。
\begin{equation}\label{b2:eq:07-jacobian}
 J(t;s,a)=\exp\left(\int_s^t\operatorname{div}_xb(X(r;s,a),r)\,\mathrm dr\right)>0.
\end{equation}
\end{proposition}
\begin{proof}
令 \(A=D_aX\)、\(B(t)=D_xb(X(t;s,a),t)\)，则可微依赖给出 \(A'=BA\)、\(A(s)=I\)。沿固定轨道，\(B\) 是连续矩阵函数，故可直接应用\cref{b2:prop:ode-liouville}。注意 \(\operatorname{tr}B(t)=\operatorname{div}_xb(X(t;s,a),t)\)，代入该命题便得到\eqref{b2:eq:07-jacobian}，并由实指数的正性得到 \(J>0\)。
\end{proof}

这就是 Liouville 公式在流的变分方程上的具体形式。散度描述局部体积变化率；散度为零的流保持 Lebesgue 体积。密度 \(\rho\) 要抵消体积变化，因此守恒质量的方程应写成\term[lianxuxing-fangcheng]{连续性方程}[continuity equation]
\begin{equation}\label{b2:eq:07-continuity}
 \rho_t+\operatorname{div}_x(b\rho)=0,
 \quad\text{即}\quad
 \rho_t+b\cdot\nabla\rho+(\operatorname{div}b)\rho=0.
\end{equation}

\begin{theorem}\label{b2:thm:07-mass-transport}
设 \(d\) 为正整数，\(T>0\)，\(b\in C^2(\mathbb R^d\times[0,T];\mathbb R^d)\)，\(|b(x,t)|\leq C(1+|x|)\) 对某个 \(C\geq0\) 成立；\(X(t;0,a)\) 为此速度场的轨道，\(J(t;0,a)=\det D_aX(t;0,a)\)。设 \(\rho_0\in C^1(\mathbb R^d;\mathbb C)\)。则初值问题 \(\rho_t+\operatorname{div}_x(b\rho)=0\)、\(\rho(\cdot,0)=\rho_0\) 有唯一的经典解 \(\rho\in C^1(\mathbb R^d\times[0,T];\mathbb C)\)，并满足
\begin{equation}\label{b2:eq:07-density-jacobian}
 \rho(X(t;0,a),t)\cdot J(t;0,a)=\rho_0(a).
\end{equation}
若 \(\rho_0\) 为非负实值函数，或 \(\rho_0\in L^1(\mathbb R^d;\mathbb C)\)，则对每个 Borel 集 \(E\subseteq\mathbb R^d\) 和 \(t\in[0,T]\)，
\[
 \int_{X(t;0,E)}\rho(x,t)\,\mathrm dx=\int_E\rho_0(a)\,\mathrm da.
\]
非负情形允许两边为 \(+\infty\)。可积情形还有 \(\|\rho(\cdot,t)\|_1=\|\rho_0\|_1\)。
\end{theorem}
\begin{proof}
因为 \(\operatorname{div}b\in C^1\)，可对\eqref{b2:eq:07-continuity}使用\cref{b2:thm:07-variable-transport}，得到经典解的存在唯一性。沿轨道，\(\rho\) 满足的标量方程与 \(J\) 的方程恰有相反的零阶系数，乘积求导为零，得到\eqref{b2:eq:07-density-jacobian}。对微分同胚 \(a\mapsto X(t;0,a)\) 使用\secref{b2:sec:0102}的换元公式，因为 \(J>0\)，得到积分等式。若初始密度非负，可直接用非负换元；若绝对可积，先对绝对值换元，得到可积性及范数恒等式，再对正、负部或实、虚部分别积分。
\end{proof}

这里加 \(C^2\) 条件是为了由显式公式得到 \(C^1\) 密度；\cref{b2:prop:07-jacobian}本身只需要 \(b\in C^1\)。对于速度 \(b(x)=x\)，输运标量为 \(u(x,t)=g(e^{-t}x)\)，而守恒密度为 \(\rho(x,t)=e^{-dt}\cdot\rho_0(e^{-t}x)\)。两者的支集作同样的扩张，后者的振幅同时降低。运动几何、所携带的值与密度三者至此有了各自准确的表达式。

\ProseParagraphBreak

\Cref{b2:tab:07-scalar-density}把没有源项时的两种未知量并列比较。表中 \(d\) 为正整数，\(T>0\)，同一实速度场 \(b\in C^2(\mathbb R^d\times[0,T];\mathbb R^d)\) 满足 \(|b(x,t)|\leq C(1+|x|)\)、\(C\geq0\)；记 \(X(t,a)=X(t;0,a)\)、\(J(t,a)=\det D_aX(t;0,a)>0\)。标量初值取 \(g\in C^1(\mathbb R^d;\mathbb C)\)，密度初值取非负的 \(\rho_0\in C^1(\mathbb R^d;\mathbb R)\cap L^1(\mathbb R^d;\mathbb R)\)。表中的“移动区域”是同一群初始粒子 \(E\) 的像 \(X(t,E)\)，其中 \(E\subseteq\mathbb R^d\) 为 Borel 集。

\begin{table}[htbp]
\centering
\small
\renewcommand{\arraystretch}{1.3}
\caption[输运标量与守恒密度]{输运标量与守恒密度。两列使用同一个流；密度沿轨道的变化抵消了体积因子，标量沿轨道的数值不变则没有这种补偿。}
\label{b2:tab:07-scalar-density}
\begin{tabularx}{\linewidth}{@{}>{\raggedright\arraybackslash}p{1.5cm}>{\raggedright\arraybackslash}X>{\raggedright\arraybackslash}X@{}}
\toprule
 & 输运标量 \(u\) & 守恒密度 \(\rho\) \\
\midrule
方程 & \(u_t+b\cdot\nabla u=0\) & \(\rho_t+\operatorname{div}_x(b\rho)=0\) \\
沿轨道 & \(u(X(t,a),t)=g(a)\) & \(\rho(X(t,a),t)J(t,a)=\rho_0(a)\) \\
\(b(x)=x\) & \(u(x,t)=g(e^{-t}x)\) & \(\rho(x,t)=e^{-dt}\rho_0(e^{-t}x)\) \\
保持的量 & 每个粒子携带的函数值；总积分或 \(L^p\) 范数即使有限，也不必保持。 & 随流移动区域内的质量：\(\int_{X(t,E)}\rho=\int_E\rho_0\)。总质量及 \(L^1\) 范数不变。 \\
\bottomrule
\end{tabularx}
\end{table}

移动区域的质量不变，并不表示每个固定观察区域内的质量也不变；粒子仍可穿过固定区域的边界。这一区别不要求移动区域具有光滑边界，质量恒等式已经由换元对所有 Borel 集给出。

\ProseParagraphBreak

Jacobi 公式同时决定各个积分范数的演化。以下取正整数 \(d\)、\(T>0\)，实速度场 \(b\in C^1(\mathbb R^d\times[0,T];\mathbb R^d)\) 满足 \(|b(x,t)|\leq C(1+|x|)\)、\(C\geq0\)，且 \(c\in C^1(\mathbb R^d\times[0,T];\mathbb R)\)、\(g\in C^1(\mathbb R^d;\mathbb C)\)。对无源问题 \(u_t+b\cdot\nabla u+cu=0\)、\(u(x,0)=g(x)\) 的经典解，沿轨道有
\[
 u(X(t;0,a),t)=g(a)\exp\left(-\int_0^tc(X(r;0,a),r)\,\mathrm{d}r\right).
\]
对 \(1\leq p<\infty\) 作换元，得到精确式
\[
 \|u(\cdot,t)\|_p^p
 =\int |g(a)|^p\exp\left(\int_0^t
   (\operatorname{div}b-pc)(X(r;0,a),r)\,\mathrm{d}r\right)\,\mathrm{d}a.
\]
若存在实常数 \(D,c_0\)，使整个时空条带上 \(\operatorname{div}b\leq D\)、\(c\geq c_0\)，且 \(g\in L^p(\mathbb R^d;\mathbb C)\)，便有 \(\|u(\cdot,t)\|_p\leq e^{(D/p-c_0)t}\|g\|_p\)。无穷范数不读取体积；若 \(g\in L^\infty\)，直接由表示公式得 \(\|u(\cdot,t)\|_\infty\leq e^{-c_0t}\|g\|_\infty\)。保持同一组 \(b,c\)，在有源情形取 \(f\in C^1(\mathbb R^d\times[0,T];\mathbb C)\)，对每个源时刻用同一换元，再用 Minkowski 不等式，给出
\[
 \|u(\cdot,t)\|_p
 \leq e^{(D/p-c_0)t}\|g\|_p
 +\int_0^t e^{(D/p-c_0)(t-r)}\|f(\cdot,r)\|_p\,\mathrm{d}r,
\]
其中 \(p=\infty\) 按 \(D/p=0\) 解释，并要求右侧有限。

\ProseParagraphBreak

对\cref{b2:thm:07-mass-transport}中 \(b\in C^2\) 所给出的守恒密度，零阶系数是 \(c=\operatorname{div}b\)。它的 \(p\) 次质量因子变为 \((1-p)\operatorname{div}b\)。所以只有 \(p=1\) 对任意体积变化恰好抵消；若 \(1\leq p\leq\infty\)、\(\rho_0\in L^p(\mathbb R^d;\mathbb C)\)，且整个时空条带上 \(|\operatorname{div}b|\leq D\)、\(D\geq0\)，则
\[
 \|\rho(\cdot,t)\|_p\leq e^{(1-1/p)Dt}\|\rho_0\|_p.
\]
对于 \(b(x)=x\)，还可直接算出 \(\|\rho(\cdot,t)\|_p=e^{-d(1-1/p)t}\|\rho_0\|_p\)，而输运标量的范数为 \(e^{dt/p}\|g\|_p\)。同一流对不同物理量产生不同范数变化，不能把所有输运都当成保持每个 \(L^p\) 范数的平移。

% !TeX root = ../../Book2.tex
\OptionalSection{一阶非线性方程}{b2:sec:0704}

在线性输运中，先由已知速度场确定轨道，再沿轨道计算数值。若速度依赖未知函数，这两步便相互耦合。沿特征线同时求解位置与函数值，可以构造局部光滑解；显式例子则会显示梯度怎样在有限时间内失去控制。

\subsection{拟线性方程}
\label{b2:subsec:070401}

考虑已在\secref{b2:sec:0104}定义的拟线性类型
\[
 u_t+b(x,t,u)\cdot\nabla u=f(x,t,u),\quad u(x,0)=g(x).
\]
这里 \(d\geq1\) 为整数，未知函数及数据 \(g\) 取实值，速度 \(b\) 取值于 \(\mathbb R^d\)，源项 \(f\) 取实值。设 \(b,f\) 在相关变量的开集上为 \(C^1\)，\(g\in C^1\)。把未知的函数值也看成一个坐标，解联立系统
\begin{equation}\label{b2:eq:07-quasilinear-system}
 \begin{aligned}
 X'&=b(X,t,Z),&X(0,a)&=a,\\
 Z'&=f(X,t,Z),&Z(0,a)&=g(a).
 \end{aligned}
\end{equation}
将时间也作为坐标，补上 \(t'=1\)，便可对向量 \((X,t,Z)\) 应用\cref{b2:thm:07-flow-derivative}，得到其对 \((t,a)\) 的 \(C^1\) 依赖。时刻零有 \(D_aX=I_d\)，故映射 \((a,t)\mapsto(X(t,a),t)\) 在每个 \((a_0,0)\) 附近的导数可逆。

\begin{proposition}\label{b2:prop:07-quasilinear-local}
设 \(d\geq1\) 为整数，\(V\subset\mathbb R^d\) 为开集，\(g\in C^1(V;\mathbb R)\)，\(a_0\in V\)，\(b,f\) 在包含 \((a_0,0,g(a_0))\) 的 \(\mathbb R^d\times\mathbb R\times\mathbb R\) 开集上属于 \(C^1\)，值域分别为 \(\mathbb R^d\)、\(\mathbb R\)。则在 \((a_0,0)\) 的某个时空邻域中，初值问题 \(u_t+b(x,t,u)\cdot\nabla u=f(x,t,u)\)、\(u(x,0)=g(x)\) 有实值 \(C^1\) 解，且该解局部唯一；任何满足相同初值的另一实值 \(C^1\) 解，在适当缩小的共同邻域内与它相等。
\end{proposition}
\begin{proof}
由\cref{b2:thm:01-inverse}，缩小邻域后可反解 \(a=A(x,t)\)。令 \(u(x,t)=Z(t,A(x,t))\)，则
\(u(X(t,a),t)=Z(t,a)\)。对此式求时间导数并代入\eqref{b2:eq:07-quasilinear-system}，得到偏微分方程。初始时 \(A(x,0)=x\)，所以 \(u(x,0)=g(x)\)。若另有经典解 \(v\)，沿满足 \(X'=b(X,t,v(X,t))\) 的曲线令 \(Z=v(X,t)\)，链式法则表明 \((X,Z)\) 满足同一联立系统和同一初值。常微分方程唯一性使二者与构造的 \((X,Z)\) 相同，再由局部反演得到 \(v=u\)。
\end{proof}

这一结论的邻域可以依赖 \(a_0\)。即使联立系统一直有解，其空间投影 \(a\mapsto X(t,a)\) 也可能先失去可逆性。下面的例子把这种障碍显示为一个分母趋零的计算。

\subsection{特征线交叉}
\label{b2:subsec:070402}

\term[burgers-fangcheng]{无黏 Burgers 方程}[inviscid Burgers equation]为
\[
 u_t+u u_x=0.
\]
取初值 \(u(x,0)=g(x)\)，其中 \(g\in C^1(\mathbb R;\mathbb R)\) 有界，考察 \(t\geq0\)。沿特征线 \(Z'=0\)，所以 \(Z(t,a)=g(a)\)、\(X(t,a)=a+tg(a)\)。在满足 \(1+tg'(a)>0\) 对所有 \(a\in\mathbb R\) 成立的时间内，\(X(t,\cdot)\) 严格递增，又因 \(|X(t,a)-a|\leq t\|g\|_\infty\) 而满射。它是全实轴的 \(C^1\) 微分同胚，局部逆随时空变量 \(C^1\) 变化，解由 \(u(X(t,a),t)=g(a)\) 给出，并满足
\begin{equation}\label{b2:eq:07-burgers-gradient}
 u_x(X(t,a),t)=\frac{g'(a)}{1+tg'(a)}.
\end{equation}
若 \(g'\geq0\)，上述构造对所有正时间有效；若 \(g'\) 的负最小值 \(m\) 在某点取到，则在 \(T_*=-1/m\) 之前有效，而该点轨道上的梯度在 \(t\uparrow T_*\) 时趋于负无穷。

\ProseParagraphBreak

具体取 \(g(a)=-\tanh a\)，有 \(m=-1\)、\(T_*=1\)。在 \(0\leq t<1\) 时，函数值始终处于 \((-1,1)\)，但原点处梯度为 \(-1/(1-t)\)。失效来自投影的导数 \(\partial_aX(t,0)=1-t\) 趋于零，即相邻初始标签之间的无穷小空间间隔被压缩。在 \(t=1\) 时，\(a\mapsto a-\tanh a\) 仍严格递增，但在原点的导数为零，逆映射已不再为 \(C^1\)。因此首次梯度爆破并不要求两条不同标签的投影恰在这一时刻相交。当 \(t>1\) 时，映射 \(a-t\tanh a\) 在原点附近下降，在两端上升，才失去单射性。两条不同初始标签的轨道可以投影到同一个时空点，而所携带的值不同，不可能成为单值经典解。\Cref{b2:fig:07-characteristic-compression}区分了光滑阶段与失效后的轨道投影。

\ProseParagraphBreak

这里交叉的是 \((x,t)\) 平面中的投影。在扩充状态 \((X,Z)\) 中，特征系统为 \(X'=Z\)、\(Z'=0\)，其流是 \((a,z_0)\mapsto(a+tz_0,z_0)\)，逆为 \((x,z)\mapsto(x-tz,z)\)。所以完整状态的流始终单射，常微分方程的唯一性并未失效；失效的是丢去 \(Z\) 坐标后，用位置反解初始标签的步骤。
% 自绘：g(a)=-tanh(a) 的特征投影；虚线仅表示光滑失效后的轨道投影。
\begin{figure}[htbp]
\centering
\begin{tikzpicture}[x=2.3cm,y=2.2cm,>=stealth]
 \draw[->] (-1.35,0)--(1.35,0) node[right]{\(x\)};
 \draw[->] (-1.3,-0.05)--(-1.3,1.6) node[above]{\(t\)};
 \draw[gray,densely dotted] (-1.25,1)--(1.25,1);
 \node[left] at (-1.3,1){\(1\)};
 \foreach \a/\v in {-1/-0.761594,-0.75/-0.635149,-0.5/-0.462117,-0.25/-0.244919,0/0,0.25/0.244919,0.5/0.462117,0.75/0.635149,1/0.761594}{
   \draw[AnalysisBlue,thick] (\a,0)--({\a-\v},1);
   \draw[AnalysisBlue,dashed] ({\a-\v},1)--({\a-1.45*\v},1.45);
 }
 \foreach \x in {-1,0,1}{\draw (\x,0.025)--(\x,-0.025) node[below]{\(\x\)};}
 \node[anchor=west,align=left] at (0.68,1.28){\small 轨道投影};
 \node[anchor=west] at (0.68,0.72){\small 光滑阶段};
\end{tikzpicture}
\caption[特征线的压缩与投影交叉]{特征线的压缩与投影交叉。初值为 \(g(a)=-\tanh a\)，轨道投影为 \(x=a-t\tanh a\)。在 \(t=1\) 首次发生梯度爆破；虚线延长只展示投影的交叉，不表示单值经典解的延续。}
\label{b2:fig:07-characteristic-compression}
\end{figure}


\subsection{激波形成的预告}
\label{b2:subsec:070403}

在经典解存在时，Burgers 方程等价于
\(u_t+\partial_x(u^2/2)=0\)。对于有跳跃的函数，后一种写法仍可通过分布导数解释。设一般通量 \(F\in C^1(\mathbb R;\mathbb R)\)，\(I\subset\mathbb R\) 为开时间区间，\(s\in C^1(I;\mathbb R)\)。令实值函数在 \(x<s(t)\) 一侧恒为常数 \(u_-\)，在 \(x>s(t)\) 一侧恒为常数 \(u_+\)；曲线上的取值不影响分布。对 \(\varphi\in C_c^\infty(\mathbb R\times I)\)，把分布 \(u_t+\partial_xF(u)\) 的测试积分在曲线两侧分别分部积分，得到界面项
\[
 \int \bigl(F(u_+)-F(u_-)-s'(t)(u_+-u_-)\bigr)
 \varphi(s(t),t)\,\mathrm dt.
\]
它对所有紧支集光滑测试函数为零，当且仅当
\[
 s'(t)(u_+-u_-)=F(u_+)-F(u_-).
\]
这称为\term[rankine-hugoniot]{Rankine--Hugoniot 条件}[Rankine--Hugoniot condition]。对不同的 Burgers 左右状态，速度为 \((u_-+u_+)/2\)；这样的移动跳跃称为\term[jibo]{激波}[shock wave]。

\ProseParagraphBreak

界面恒等式只给出弱方程的相容条件，没有证明某个跳跃是光滑解失效后的唯一延续。物理上允许的压缩跳跃与不允许的展开跳跃，都可能满足同一恒等式。熵条件用来区分这些跳跃；\secref{b2:sec:0705}将从两常数初值出发，计算激波与稀疏波，并解释相应的选择。

\subsection{Hamilton--Jacobi 方程}
\label{b2:subsec:070404}

另一类重要的一阶非线性方程是\term[hamilton-jacobi-fangcheng]{Hamilton--Jacobi 方程}[Hamilton--Jacobi equation]
\[
 u_t+H(x,\nabla u)=0.
\]
这里 \(d\geq1\) 为整数，\(H\) 与未知函数 \(u\) 都取实值。设 \(V\subset\mathbb R^d\) 开、\(a_0\in V\)，\(g\in C^3(V;\mathbb R)\)，\(H\) 是包含 \((a_0,\nabla g(a_0))\) 的 \(\mathbb R^d\times\mathbb R^d\) 开邻域上的实值 \(C^3\) 函数，只在初始点附近构造光滑解。取
\begin{equation}\label{b2:eq:07-hamilton-system}
 X'=H_p(X,P),\quad P'=-H_x(X,P),\quad
 Z'=P\cdot H_p(X,P)-H(X,P),
\end{equation}
初始值为 \(X(0,a)=a\)、\(P(0,a)=\nabla g(a)\)、\(Z(0,a)=g(a)\)。前两式描述位置和候选梯度的演化，第三式计算函数值。

\begin{proposition}\label{b2:prop:07-hamilton-local}
设 \(d\geq1\) 为整数，\(V\subset\mathbb R^d\) 为开集，\(g\in C^3(V;\mathbb R)\)，\(a_0\in V\)。设 \(\Omega\subset\mathbb R^d\times\mathbb R^d\) 为包含 \((a_0,\nabla g(a_0))\) 的开集，\(H\in C^3(\Omega;\mathbb R)\)。取 \(X(0,a)=a\)、\(P(0,a)=\nabla g(a)\)、\(Z(0,a)=g(a)\)，系统\eqref{b2:eq:07-hamilton-system}在 \((a_0,0)\) 的某个时空邻域中确定一个实值 \(C^2\) 解 \(u\)，满足 \(u_t+H(x,\nabla u)=0\)、\(u(x,0)=g(x)\) 及 \((x,\nabla u(x,t))\in\Omega\)。该解在实值 \(C^2\) 类中局部唯一：若另一实值 \(C^2\) 解 \(v\) 定义在 \((a_0,0)\) 附近，满足同一方程和初值，且 \((x,\nabla v(x,t))\in\Omega\)，则 \(v\) 在适当缩小的共同邻域内与 \(u\) 相等。
\end{proposition}
\begin{proof}
先核对特征系统的二阶参数依赖。令 \(N=2d+1\)，把独立初始状态写成 \(y_0=(x_0,p_0,z_0)\in\Omega\times\mathbb R\)，并设
\[
 B(x,p,z)=\bigl(H_p(x,p),-H_x(x,p),p\cdot H_p(x,p)-H(x,p)\bigr).
\]
由 \(H\in C^3\)，向量场 \(B\) 属于 \(C^2\)。记 \(\Phi_t(y_0)\) 为 \(Y'=B(Y)\)、\(Y(0)=y_0\) 的局部流。由\cref{b2:thm:07-flow-derivative}，流为 \(C^1\)，且 \(L(t,y_0)=D_{y_0}\Phi_t(y_0)\) 满足
\[
 L'=DB(\Phi_t(y_0))L,\quad L(0,y_0)=I_N.
\]
对独立的方向 \(r,s\in\mathbb R^N\)，二阶导数的候选值 \(V_{r,s}\) 由下列连续系数线性积分方程的等价微分形式确定；其存在唯一性用一阶变分方程相同的积分迭代得到：
\[
 V_{r,s}'=DB(\Phi_t(y_0))V_{r,s}
       +D^2B(\Phi_t(y_0))[Lr,Ls],
 \quad V_{r,s}(0)=0.
\]
零初值的原因是 \(\Phi_0(y_0)=y_0\) 对独立变量 \(y_0\) 为恒等映射。

为验证候选值确实是二阶导数，缩小初始邻域与时间区间，使所有流及其连接线段位于同一个紧球 \(K\subset\Omega\times\mathbb R\) 中。记 \(Y_\eta(t)=\Phi_t(y_0+\eta s)\)、\(Y(t)=\Phi_t(y_0)\)，其中 \(\eta\) 为趋于零的非零实数。一阶 \(C^1\) 依赖给出 \((Y_\eta-Y)/\eta\to Ls\)，在该紧时间区间上为一致收敛。对一阶变分方程作差，令
\[
 Q_\eta(t)=\frac{L(t,y_0+\eta s)r-L(t,y_0)r}{\eta},
\]
则
\[
 Q_\eta'=DB(Y_\eta)Q_\eta
       +\frac{DB(Y_\eta)-DB(Y)}{\eta}Lr,
 \quad Q_\eta(0)=0.
\]
中值积分将这里的算子差商写成
\[
 \frac{DB(Y_\eta)-DB(Y)}{\eta}
 =\int_0^1D^2B\bigl(Y+\theta(Y_\eta-Y)\bigr)
       \left[\frac{Y_\eta-Y}{\eta},\,\cdot\right]\,\mathrm{d}\theta.
\]
\(D^2B\) 在 \(K\) 上一致连续，所以该算子差商一致趋于 \(D^2B(Y)[Ls,\cdot]\)；同时 \(DB(Y_\eta)\to DB(Y)\) 一致。二阶微分的对称性使极限非齐次项 \(D^2B(Y)[Ls,Lr]\) 正是 \(V_{r,s}\) 方程中的那一项。把 \(Q_\eta\) 的积分方程与 \(V_{r,s}\) 的方程相减，所有系数及非齐次项均有共同上界，余项一致趋零，Gronwall 估计遂给出 \(Q_\eta\to V_{r,s}\) 一致。方程中系数与初始状态连续，作差后再用同一估计，得到 \(V_{r,s}\) 对 \((t,y_0,r,s)\) 连续。因此 \(D_{y_0}\Phi_t\) 的各方向偏导数连续，\(\Phi_t\) 对 \(y_0\) 为 \(C^2\)。再由 \(\partial_t\Phi_t=B(\Phi_t)\) 和一阶变分方程，混合二阶导数与二阶时间导数也连续，故在该共同邻域内，流对 \((t,y_0)\) 联合为 \(C^2\)。

现在才代入实际初始映射。先由连续性缩小 \(a_0\) 的邻域，使 \((a,\nabla g(a))\) 均落在 \(\Omega\) 内，并使下面的初始状态落在上述局部流的共同初始邻域中。令
\[
 y_0(a)=(a,\nabla g(a),g(a)),\quad
 Y(t,a)=\Phi_t(y_0(a)).
\]
由于 \(g\in C^3\)，映射 \(y_0\) 为 \(C^2\)。二阶链式法则给出
\[
 \begin{aligned}
 D_a^2Y(t,a)[v,w]
 &=D_{y_0}^2\Phi_t(y_0(a))[Dy_0(a)v,Dy_0(a)w]\\
 &\quad+D_{y_0}\Phi_t(y_0(a))D^2y_0(a)[v,w],
 \quad v,w\in\mathbb R^d.
 \end{aligned}
\]
特别地，初始时刻的结果是
\[
 D_a^2Y(0,a)[v,w]
 =\bigl(0,D^3g(a)[v,w],D^2g(a)[v,w]\bigr).
\]
这里 \(D^3g(a)[v,w]\) 表示向量 \(D^2(\nabla g)(a)[v,w]\)，其第 \(i\) 个分量为 \(D^2(\partial_i g)(a)[v,w]\)。这个初值一般不为零；为零的是独立状态流的 \(D_{y_0}^2\Phi_0\)。这一复合论证证明了 \(X,P,Z\) 对 \((t,a)\) 为 \(C^2\)。

从 \(D_aX(0,a)=I_d\) 出发，对 \((a,t)\mapsto(X(t,a),t)\) 应用\cref{b2:thm:01-inverse}，得到局部逆 \((x,t)\mapsto(A(x,t),t)\)。原映射为 \(C^2\)，逆微分公式说明逆映射的导数为 \(C^1\)，故逆也为 \(C^2\)。令 \(u(x,t)=Z(t,A(x,t))\)，便有 \(u\in C^2\) 及 \(u(x,0)=g(x)\)。关键是验证 \(P=\nabla u\)，不能把候选梯度直接当成真实梯度。对每个初始坐标 \(a_j\)，记
\(W_j=\partial_{a_j}Z-P\cdot\partial_{a_j}X\)。对系统求导，利用 \(X'=H_p\)、\(P'=-H_x\)，得到
\[
 \begin{aligned}
 \partial_{a_j}Z'
 &=P\cdot\partial_{a_j}(H_p)-H_x\cdot\partial_{a_j}X,\\
 (P\cdot\partial_{a_j}X)'
 &=-H_x\cdot\partial_{a_j}X+P\cdot\partial_{a_j}(H_p).
 \end{aligned}
\]
故 \(W_j'=0\)，初值又为 \(\partial_jg-\partial_jg=0\)。因此 \(D_aZ=P^{\mathsf T}D_aX\)，由 \(D_aX\) 可逆得到 \(\nabla_xu(X,t)=P\)。最后沿轨道有
\(u_t+\nabla u\cdot X'=Z'=P\cdot H_p-H\)，约去相同项得 \(u_t+H(x,\nabla u)=0\)。

最后证明局部唯一性。设 \(v\) 是在 \((a_0,0)\) 附近满足同一初值问题的另一实值 \(C^2\) 解，且 \((x,\nabla v(x,t))\in\Omega\)。对它的方程求空间梯度，得到
\[
 \partial_t\nabla v+D_x^2v\,H_p(x,\nabla v)+H_x(x,\nabla v)=0.
\]
速度 \((x,t)\mapsto H_p(x,\nabla v(x,t))\) 为 \(C^1\)，故从每个邻近标签 \(a\) 出发可以解
\[
 \widetilde X'=H_p(\widetilde X,\nabla v(\widetilde X,t)),
 \quad \widetilde X(0,a)=a.
\]
缩小邻域使这些短轨道始终留在 \(v\) 的定义域内，并令 \(\widetilde P(t,a)=\nabla v(\widetilde X(t,a),t)\)、\(\widetilde Z(t,a)=v(\widetilde X(t,a),t)\)。链式法则与梯度方程分别给出
\[
 \widetilde P'=-H_x(\widetilde X,\widetilde P),\quad
 \widetilde Z'=\widetilde P\cdot H_p(\widetilde X,\widetilde P)
                    -H(\widetilde X,\widetilde P).
\]
其初始状态同为 \((a,\nabla g(a),g(a))\)。向量场 \(B\) 的常微分方程唯一性使 \((\widetilde X,\widetilde P,\widetilde Z)=(X,P,Z)\)。因此 \(v(X(t,a),t)=Z(t,a)=u(X(t,a),t)\)，局部投影的逆坐标又覆盖了一个时空邻域，遂得到该邻域内 \(v=u\)。
\end{proof}

取 \(H(p)=|p|^2/2\) 时，\(P=\nabla g(a)\)、\(X=a+t\nabla g(a)\)、\(Z=g(a)+t|\nabla g(a)|^2/2\)。可逆性要求 \(I+tD^2g(a)\) 不退化。这与 Burgers 方程的分母具有同一几何含义，但高维中须检查整个矩阵。对于标量守恒律，下一节将用弱解与熵条件研究特征线交叉之后的延续。

\subsection{守恒律的位置}
\label{b2:subsec:070405}

设 \(d\geq1\) 为整数，\(F\in C^1(\mathbb R;\mathbb R^d)\)。形如 \(u_t+\operatorname{div}_xF(u)=0\) 的方程称为\term[shouhenglv]{守恒律}[conservation law]。若在时空开集 \(\Omega\subset\mathbb R^d\times\mathbb R\) 上 \(u\) 为实值 \(C^1\) 函数，则由链式法则它等价于拟线性方程 \(u_t+F'(u)\cdot\nabla u=0\)；不光滑时，两种写法的意义不同，后一项中的乘积未必已有定义。守恒形式可直接作用于测试函数，通过积分表达局部通量平衡。

\ProseParagraphBreak

线性连续性方程的速度由外部给定，密度与体积的补偿由 Jacobi 公式控制；非线性守恒律的速度随状态变化，投影交叉会使经典表示失效。两者都使用特征线，却不能把线性流的全局微分同胚结论原样移到非线性问题。本节建立了光滑阶段的局部构造和一个明确的失效机制，为进一步学习守恒律与 Hamilton--Jacobi 方程作准备。

% !TeX root = ../../Book2.tex
\OptionalSection{熵解、激波与 Riemann 问题}{b2:sec:0705}

待撰写。

\subsection{守恒律弱解与非唯一性}
\label{b2:subsec:070501}

待撰写。

\subsection{Burgers 方程的 Riemann 问题}
\label{b2:subsec:070502}

待撰写。

\subsection{熵条件与激波选择}
\label{b2:subsec:070503}

待撰写。

\subsection{一般标量守恒律的理论边界}
\label{b2:subsec:070504}

待撰写。


\begin{chaptersummary}
特征线法将一阶偏微分方程拉回常微分方程。求出轨道只是第一步：还要确认每个点能沿轨道返回数据面，轨道之间由唯一性保证不分叉，流对初始位置的导数又允许把沿线公式还原成经典解。\chapref{b2:ch:ode}给出这些流的基本性质；本章把它们用于表示公式与非特征数据问题。局部存在不能直接保证全局回溯，增长条件仍须单独核对。

\ProseParagraphBreak

输运标量在齐次情形沿轨道保持数值，守恒密度则需乘上流的 Jacobi 行列式才保持不变。源项不断产生新数值，零阶项沿历史轨道累积为指数因子。支集由流携带，统一速度界才给出统一传播圆锥。在选读的\secref{b2:sec:0704}中，非线性方程的空间投影还可能失去可逆性，光滑解的构造因此具有明确的时间边界；略过该节不影响前面的线性流与质量输运结论。
\end{chaptersummary}

\BookExercises

前两题训练边界数据与变换语言的接合；中间四题分别检验全局条件、体积变化、源项及闭轨道相容性。最后两题使用选读节，比较光滑失效与弱方程的可能延续。各题均在本章已建立的工具内给出解答。

% 自拟：把全空间常速度问题改为半直线入流问题，检查角点一阶相容。
\Needspace{6\baselineskip}
\begin{exercise}\label{b2:ex:07-inflow}
设 \(c>0\)，在 \(x,t\geq0\) 上求解 \(u_t+c u_x=0\)，初值为 \(u(x,0)=g(x)\)，边界值为 \(u(0,t)=h(t)\)，其中 \(g,h\in C^1([0,\infty);\mathbb C)\)。给出分片公式，并证明它为闭象限上的 \(C^1\) 解的充要条件是
\(g(0)=h(0)\)、\(h'(0)=-c g'(0)\)。
\end{exercise}
\begin{solution}
回溯直线先遇到 \(t=0\) 或 \(x=0\)，故
\[
 u(x,t)=\begin{cases}g(x-ct),&x\geq ct,\\h(t-x/c),&x\leq ct.\end{cases}
\]
在接合线 \(x=ct\) 上，连续性等价于 \(g(0)=h(0)\)。两侧梯度分别趋于 \((g'(0),-c g'(0))\) 与 \((-h'(0)/c,h'(0))\)，相等恰好是第二条件。满足这两项时，沿接合线的两侧一阶导数连续接合，函数为 \(C^1\)，并逐点满足方程。任何经典解沿轨道恒定，所以分片公式也证明唯一性。
\end{solution}

% 自拟：将特征平移与第二编的单位周期 Fourier 归一化连接。
\Needspace{6\baselineskip}
\begin{exercise}\label{b2:ex:07-fourier-translation}
设 \(\mathbb T=\mathbb R/\mathbb Z\) 配备总质量为 \(1\) 的 Lebesgue 测度，\(g\in C^1(\mathbb T;\mathbb C)\)，\(c\in\mathbb R\)。证明周期输运解 \(u(x,t)=g(x-ct)\) 对每个 \(k\in\mathbb Z\)、\(t\in\mathbb R\) 的 Fourier 系数为
\(\widehat u(k,t)=e^{-2\pi\mathrm i kct}\widehat g(k)\)。进而说明对每个非负整数 \(N\)，Fourier 部分和 \(S_Nu\) 的误差在每个时刻具有与初始时刻相同的 \(L^2(\mathbb T)\) 范数。
\end{exercise}
\begin{solution}
在 \(\int_0^1g(x-ct)e^{-2\pi\mathrm i kx}\,\mathrm dx\) 中令 \(y=x-ct\)。由于 \(g(y)e^{-2\pi\mathrm i ky}\) 为单位周期，换元后的任意长度为一的积分都等于 \(\widehat g(k)\)，而剩余相位为 \(e^{-2\pi\mathrm i kct}\)。所以 \(S_Nu(x,t)=S_Ng(x-ct)\)。对二者之差再用圆周平移不变性，得到 \(\|u(\cdot,t)-S_Nu(\cdot,t)\|_2=\|g-S_Ng\|_2\)。这里不需要逐项微分一个无限级数。
\end{solution}

% 自拟：由有限时间逃逸构造零初值非唯一，展示全局增长假设的真实用途。
\Needspace{6\baselineskip}
\begin{exercise}\label{b2:ex:07-infinity-inflow}
在 \(\mathbb R\times[0,\infty)\) 上考虑实值方程 \(u_t+(1+x^2)u_x=0\)、\(u(x,0)=0\)。构造非零的 \(u\in C^1(\mathbb R\times[0,\infty);\mathbb R)\)，并解释为何不违反本章的全局唯一性定理。
\end{exercise}
\begin{solution}
取 \(F\in C^\infty(\mathbb R)\)，令 \(F(s)=0\) 当 \(s\leq\pi/2\)，令 \(F(s)=\exp(-(s-\pi/2)^{-2})\) 当 \(s>\pi/2\)。定义 \(u(x,t)=F(t-\arctan x)\)，则 \(u_t=F'\)、\(u_x=-F'/(1+x^2)\)，方程成立。初始时 \(-\arctan x<\pi/2\)，故初值为零；取 \(t>\pi/2\)、\(x=0\) 则非零。速度为二次增长，轨道 \(\arctan X=t+\arctan a\) 在有限时间到达无穷，某些终点不能回溯到初始面。全局唯一性所需的完整回溯轨道因此缺失。
\end{solution}

% 自拟：剪切流中守恒质量与非守恒矩相区分。
\Needspace{6\baselineskip}
\begin{exercise}\label{b2:ex:07-shear-moment}
设速度场 \(b(x,y)=(y,0)\)，\(\rho_0\in C_c^1(\mathbb R^2;\mathbb R)\)。求 \(\rho_t+\operatorname{div}(b\rho)=0\)、\(\rho(\cdot,0)=\rho_0\) 在 \(t\geq0\) 的经典解，并用换元证明其总积分不变，而第一矩满足
\[
 \int_{\mathbb R^2}x\rho(x,y,t)\,\mathrm dx\,\mathrm dy
 =\int_{\mathbb R^2}(a+tb)\rho_0(a,b)\,\mathrm da\,\mathrm db.
\]
这里积分变量 \(b\) 是实数坐标，非前面的向量场。
\end{exercise}
\begin{solution}
轨道为 \(X(t;0,(a,b))=(a+tb,b)\)，其导数矩阵为 \(\left(\begin{smallmatrix}1&t\\0&1\end{smallmatrix}\right)\)，行列式等于一。故 \(\rho(x,y,t)=\rho_0(x-ty,y)\)。以 \(x=a+tb,y=b\) 换元，质量等于 \(\int\rho_0\)，而把积分中的 \(x\) 同时换为 \(a+tb\) 就得到矩公式。紧支集保证这些积分绝对收敛。散度为零只保证体积和质量不变，并不要求重心静止。
\end{solution}

% 自拟：含质量源的移动控制体，实质增加源项与Jacobi的组合。
\Needspace{6\baselineskip}
\begin{exercise}\label{b2:ex:07-moving-source}
设 \(d\) 为正整数，\(T>0\)，\(b\in C^2(\mathbb R^d\times[0,T];\mathbb R^d)\)，且对某个 \(C\geq0\) 有 \(|b(x,t)|\leq C(1+|x|)\)。记 \(X(t;0,a)\) 为从 \(a\) 出发的轨道，\(J(t;0,a)=\det D_aX(t;0,a)\)，即\cref{b2:thm:07-mass-transport}中使用的流与体积因子。设 \(\rho,q\in C^1(\mathbb R^d\times[0,T];\mathbb C)\) 满足 \(\rho_t+\operatorname{div}(b\rho)=q\)。对任意有界 Borel 集 \(E\subseteq\mathbb R^d\) 和 \(0\leq t\leq T\)，证明
\[
 \frac{\mathrm d}{\mathrm dt}\int_{X(t;0,E)}\rho(x,t)\,\mathrm dx
 =\int_{X(t;0,E)}q(x,t)\,\mathrm dx.
\]
端点处取相应单侧导数。说明此式无需 \(E\) 的边界光滑。
\end{exercise}
\begin{solution}
在固定集合 \(E\) 上换元，左边积分成为 \(\int_E\rho(X(t;0,a),t)J(t;0,a)\,\mathrm da\)。乘积求导与 Jacobi 方程给出被积函数的时间导数等于 \(q(X(t;0,a),t)J(t;0,a)\)。由于 \(\overline E\) 紧且时间限于紧区间，所有函数与导数在相应轨道集上有界，故积分号下求导合法。再换元得到结论。论证没有使用边界积分，因此 Borel 可测和有界性已足够。
\end{solution}

% 自拟：闭轨道的全局相容，超出局部特征线公式。
\Needspace{6\baselineskip}
\begin{exercise}\label{b2:ex:07-closed-orbits}
在环域 \(\Omega=\{(x,y)\in\mathbb R^2:1<\sqrt{x^2+y^2}<2\}\) 上给定 \(f\in C^1(\Omega;\mathbb C)\)。证明稳态方程 \(-yu_x+xu_y=f\) 存在单值解 \(u\in C^1(\Omega;\mathbb C)\)，当且仅当对每个 \(1<r<2\)，
\[
 \int_0^{2\pi}f(r\cos\theta,r\sin\theta)\,\mathrm d\theta=0.
\]
在存在时写出全部解。
\end{exercise}
\begin{solution}
极坐标中算子为 \(\partial_\theta\)，单值函数必须以 \(2\pi\) 为周期，故沿一周积分给出必要条件。若条件成立，令
\[
 U(r,\theta)=A(r)+\int_0^\theta f(r\cos s,r\sin s)\,\mathrm ds,
\]
其中 \(A\in C^1((1,2);\mathbb C)\) 任意。零均值保证 \(U\) 周期；对 \(r\) 的导数也周期，因为零均值恒等式可对 \(r\) 求导。故它下降为环域上的 \(C^1\) 函数并满足方程。任意两个解之差在每个圆上恒定，所以全体解恰由 \(A(r)\) 给出。
\end{solution}

% 自拟选读：同一阶跃初值的两个弱解，展示仅用跳跃关系不保证唯一。
\Needspace{6\baselineskip}
\begin{optionalexercise}\label{b2:ex:07-expansion-nonunique}
取实值 Burgers 初值 \(g(x)=0\) 当 \(x\leq0\)，\(g(x)=1\) 当 \(x>0\)。证明下列两个实值有界函数在 \(\mathbb R\times(0,\infty)\) 上都满足分布方程 \(u_t+\partial_x(u^2/2)=0\)，且当 \(t\downarrow0\) 时在 \(L^1_{\mathrm{loc}}(\mathbb R)\) 中收敛到同一初值：
\[
 u_1(x,t)=\mathbf1_{\{x>t/2\}},\quad
 u_2(x,t)=\begin{cases}0,&x\leq0,\\x/t,&0<x<t,\\1,&x\geq t.\end{cases}
\]
\end{optionalexercise}
\begin{solution}
对 \(u_1\)，左右状态为 \(0,1\)，界面速度为 \(1/2=(F(1)-F(0))/(1-0)\)，所以 Rankine--Hugoniot 条件成立；两侧方程恒成立，界面分部积分证明整个分布方程。对 \(u_2\)，扇形内部有 \(u_t=-x/t^2\)、\(uu_x=x/t^2\)，其余区域为常数。两条分界线上函数及通量连续，逐区分部积分时界面项相消，故也是弱解。测试函数的支集在 \(t>0\) 内，与 \(t=0\) 有正距离，所以逐区计算没有端点奇性。两函数与初值的差都只支撑于 \([0,t]\)，绝对值至多一，局部 \(L^1\) 误差不超过 \(t\)。二者在一个非空开集上不同，因此弱方程与该初值尚不能保证唯一。
\end{solution}

% 自拟选读：矩阵特征值控制Hamilton--Jacobi二次解的最大正时间。
\Needspace{6\baselineskip}
\begin{optionalexercise}\label{b2:ex:07-quadratic-hamilton}
设 \(d\) 为正整数，\(K\in\mathbb R^{d\times d}\) 为实对称矩阵，\(I\) 为 \(d\) 阶单位矩阵。对 \(x\in\mathbb R^d\) 上的实值方程 \(u_t+|\nabla u|^2/2=0\)、\(u(x,0)=x^{\mathsf T}Kx/2\)，在开集 \(\{t\in\mathbb R:\det(I+tK)\ne0\}\) 中含 \(0\) 的连通分支内求经典解。若 \(K\) 有负特征值，确定首个正退化时间。
\end{optionalexercise}
\begin{solution}
特征系统给出 \(X=(I+tK)a\)、\(P=Ka\)、\(Z=a^{\mathsf T}Ka/2+t|Ka|^2/2\)。代入 \(a=(I+tK)^{-1}x\)，利用 \(K\) 与 \(I+tK\) 及其逆交换，得到
\[
 u(x,t)=\frac12x^{\mathsf T}K(I+tK)^{-1}x.
\]
令 \(M=K(I+tK)^{-1}\)，则 \(M\) 对称且 \(M'=-M^2\)。故 \(u_t=-x^{\mathsf T}M^2x/2\)、\(|\nabla u|^2=x^{\mathsf T}M^2x\)，直接验证方程。若最小特征值为 \(\lambda_{\min}<0\)，首个正退化时间为 \(-1/\lambda_{\min}\)，相应方向上的 Hesse 特征值 \(\lambda_{\min}/(1+t\lambda_{\min})\) 趋于负无穷。若 \(K\geq0\)，公式对所有正时间存在。
\end{solution}
